\section{Introduction}
\label{sec:intro}

The Zinc+ PIOP, in its integer instantiation, proves the satisfaction
of an $\F_q[X]$-typed UAIR by reducing it through three transcript
projections (prime projection, evaluation projection, multipoint
evaluation) to a list of \emph{MLE evaluation claims}
$\MLE[\psia(v_g)](r^*) = a_g$, one per committed trace column $v_g$.
Those claims are discharged in the protocol's final step by the
\emph{Zip+ opening}: a Brakedown-style proximity check that lifts each
claim through the homomorphism $\Z \to \F_q$ to an integer-typed
inner product, encodes the result through the same linear code used at
commit time, and ties the inner product to the committed codeword
matrix via Merkle column openings.

\paragraph{The $\F_2$ variant.} A parallel variant of the protocol,
designed for UAIRs whose every column is typed in
$\cellring = \F_2[X]/(X^{\Dval})$, runs the analogous reduction with
$\F_q$ replaced by the binary field $\GFt = \F_2[X]/(P(X))$ where
$P(X) = X^{192} + X^7 + X^2 + X + 1$. The first six steps of the
pipeline (commit, ideal check, evaluation projection, sumcheck) carry
over almost mechanically — the ring of constants changes, the
randomness changes, the arithmetic is unchanged in structure. The
seventh step does not. The integer-side opening relies crucially on
the ring homomorphism $\Z \to \F_q$, $n \mapsto n \bmod q$, which the
proximity check uses to commute encoding through projection:
\[
  \encode\big(\,\textstyle\sum_i c_i \cdot \mathrm{row}_i \bmod q\,\big)
  \;=\;
  \big(\textstyle\sum_i c_i \cdot \encode(\mathrm{row}_i)\big) \bmod q.
\]
Translating this to $\F_2[X] \to \GFt$, $p(X) \mapsto p(\alpha)$ at a
transcript-fresh $\alpha \in \GFt$ does not yield a useful identity:
the canonical degree-$<192$ representative of a $\GFt$-element $g$
evaluates at $\alpha$ to $g$ only when $\alpha$ happens to coincide
with the field's quotient generator. For a generic $\alpha$ the bit
patterns of $g$ and of its evaluation at $\alpha$ disagree.

\paragraph{Our contribution.} We construct a Step~7 opening for the
$\F_2$ variant that:
\begin{enumerate}[leftmargin=2em]
  \item Runs entirely in $\F_2[X]$-arithmetic, without ever passing
    through $\GFt$ inside the proximity check.
  \item Lifts each per-column MLE-evaluation claim $a_g \in \GFt$
    to a single unreduced $\F_2[X]$ polynomial $a_g'$ of degree
    $< D + 2\cdot 192 - 1$. The verifier discharges the original claim
    by one local evaluation $\psia(a_g') \stackrel{?}{=} a_g$ in
    $\GFt$.
  \item Reuses the existing $\F_2$-RAA codeword matrix from commit
    time. No re-commitment, no re-encoding: the code is uniformly
    $\F_2[X]$-linear, so its action commutes through the lift.
  \item Folds all per-column claims into a single $\F_2[X]$-opening
    triple via a transcript-fresh $\GFt$-batching challenge
    $\gamma$. The communication is then $O(1)$ in the number of
    columns and linear in the row count and in the column-opening
    parameter.
\end{enumerate}
The construction's load-bearing technical content is in the lift: the
canonical bit-pattern map $\GFt \to \polyringat{192}$ does not invert
the evaluation-at-$\alpha$ map at a generic $\alpha$, so we replace it
with the unique $\F_2$-linear inverse, computed once per $\alpha$ via
a $192 \times 192$ matrix inversion over $\F_2$. The cost of this
inversion is amortised across an arbitrary number of per-entry lifts;
the per-entry lift itself is a single $192$-bit dot product. See
\Cref{sec:lift} for the analysis.

\paragraph{Why the $\F_2$ pipeline.} The motivation for an all-$\F_2$
variant is to eliminate the integer-arithmetic overhead from
proofs over computations that are themselves naturally bit-oriented
— hash functions, bit-wise circuits, ECDSA signature verification.
The integer pipeline pays a constant-factor cost in the
$\F_q[X]\to\F_q$ projection at every row of the trace; the $\F_2$
variant eliminates this projection entirely. The price is paid at
Step~7, where the bit-pattern shortcut available in the integer
pipeline does not transfer; the construction of this document is the
adapter that recovers the savings.

\paragraph{Organisation.} \Cref{sec:setting} fixes the field, ring, and
trace setup. \Cref{sec:lift} develops the $\alpha$-dependent inverse
lift $\GFt \to \polyringat{192}$ that underwrites the construction.
\Cref{sec:construction} gives the full $\gamma$-batched opening
protocol, presented as four checks plus a Merkle-bound proximity step.
\Cref{sec:soundness} sketches soundness, identifying the four failure
modes and bounding each. \Cref{sec:cost} reports communication cost
and concrete numbers for $\Dval = 32$, $\nu = 16$, rate $1/4$.
